\section{Definitions and similarity geometry}
Throughout, $H\in\C^{n\times n}$ is invertible unless explicitly noted; $H^\dagger$
denotes the conjugate transpose. All unadorned norms and condition numbers are induced
by the Euclidean vector norm. In particular,
\begin{equation}\label{eq:kappa}
 \kg(A)=\norm{A}_2\norm{A^{-1}}_2,
 \qquad A\in\C^{n\times n}\text{ invertible}.
\end{equation}

\begin{definition}
A matrix $H$ is \emph{quasi-Hermitian} with respect to a positive metric if there is an
$\eta=\eta^\dagger>0$ satisfying
\begin{equation}\label{eq:quasi}
 H^\dagger\eta=\eta H.
\end{equation}
Set $\rho=\eta^{1/2}$ and define the metric vector and operator norms by
\begin{equation}\label{eq:norms}
 \norm{x}_\eta=\sqrt{x^\dagger\eta x}=\norm{\rho x}_2,
 \qquad \norm{A}_\eta=\sup_{x\ne 0}\frac{\norm{Ax}_\eta}{\norm{x}_\eta}
 =\norm{\rho A\rho^{-1}}_2.
\end{equation}
\end{definition}

\begin{proposition}[Hermitian representative and conditioning]\label{prop:similarity}
Suppose~\eqref{eq:quasi} holds. Then
\begin{equation}\label{eq:representative}
 h=\rho H\rho^{-1}=h^\dagger,\qquad
 H^{-1}=\rho^{-1}h^{-1}\rho,
\end{equation}
and
\begin{align}
 \norm{H^{\pm1}}_\eta&=\norm{h^{\pm1}}_2,\label{eq:physicalnorm}\\
 \kg(h)&=\frac{\max_{\lambda\in\spec(H)}\abs{\lambda}}
 {\min_{\lambda\in\spec(H)}\abs{\lambda}},\label{eq:eigs}\\
 \kg(H)&\leq \kg(\eta)\kg(h),\label{eq:kappabound}\\
 \norm{H^{-1}}_2&\leq
 \frac{\sqrt{\kg(\eta)}}{\min_{\lambda\in\spec(H)}\abs{\lambda}}.
 \label{eq:inversebound}
\end{align}
\end{proposition}
\begin{proof}
Equation~\eqref{eq:quasi} and $\eta=\rho^2$ give
$H^\dagger=\rho^2H\rho^{-2}$. Thus
$h^\dagger=\rho^{-1}H^\dagger\rho=\rho H\rho^{-1}=h$.
The remaining identities follow from~\eqref{eq:norms}, the spectral theorem for $h$,
and submultiplicativity, using
$\kg(\rho)^2=\kg(\eta)$ and
$\norm{\rho^{-1}}_2\norm{\rho}_2=\sqrt{\kg(\eta)}$.
\end{proof}

\begin{remark}[Physical versus reference norm]
The metric-induced condition number of $H$ is exactly $\kg(h)$; it is
\emph{independent of which admissible metric is selected}, because any Hermitian
representative has the same real eigenvalues. Metric optimization can improve a bound
in the \emph{fixed Euclidean representation} or a specified implementation procedure,
but it cannot change these eigenvalues. Choosing $\eta$ is also constrained if
other physical observables and their inner product have already been fixed.
\end{remark}

\begin{proposition}[Admissible metrics for simple real spectrum]\label{prop:weights}
If $H=V\Lambda V^{-1}$ has distinct real eigenvalues, all positive solutions of
\eqref{eq:quasi} are
\begin{equation}\label{eq:weights}
 \eta=V^{-\dagger}WV^{-1},\qquad
 W=\diag(w_1,\ldots,w_n),\quad w_j>0.
\end{equation}
\end{proposition}
\begin{proof}
Write $G=V^\dagger\eta V$. The intertwining equation becomes
$\Lambda G=G\Lambda$. Since the eigenvalues of $\Lambda$ are distinct,
$G$ is diagonal; positivity implies its diagonal entries are strictly positive.
Conversely, every matrix in~\eqref{eq:weights} is positive and satisfies
\eqref{eq:quasi}.
\end{proof}

\begin{proposition}[A semidefinite formulation of minimum metric conditioning]
\label{prop:sdp}
Let $H$ be diagonalizable with real spectrum and define
\begin{equation}\label{eq:kappastar}
 \kappa_*(H)=\min_{\eta=\eta^\dagger>0,\;H^\dagger\eta=\eta H}
 \kappa_2(\eta).
\end{equation}
The minimum is attained and admits the equivalent characterization
\begin{equation}\label{eq:sdp}
 \kappa_*(H)=\min\left\{t\geq1:\;\exists\eta=\eta^\dagger,
 \quad H^\dagger\eta=\eta H,\quad
 \I\preceq\eta\preceq t\I\right\}.
\end{equation}
For fixed $t$, feasibility in~\eqref{eq:sdp} consists of linear matrix
inequalities and a linear matrix equation.
\end{proposition}
\begin{proof}
For any admissible positive $\eta$, division by its smallest eigenvalue
preserves quasi-Hermiticity and yields $\I\preceq\eta\preceq
\kappa_2(\eta)\I$. Conversely any feasible pair $(t,\eta)$ has
$\kappa_2(\eta)\leq t$. Existence of an admissible metric follows from
the real diagonalization. Normalize a minimizing sequence to have smallest
eigenvalue one, and choose a fixed feasible upper bound $t_0$.
The normalized metrics lie in the compact set
$\{\eta=\eta^\dagger:\I\preceq\eta\preceq t_0\I,\;
 H^\dagger\eta=\eta H\}$, so a convergent subsequence attains the
minimum by continuity of $\kappa_2$.
\end{proof}
\begin{remark}
The optimization~\eqref{eq:sdp} is a useful computable reformulation, not
a claim that selecting or optimizing a quasi-Hermitian metric is new; see
\cite{Krejcirik2018}. The condition-number objective differs from their
Hilbert--Schmidt anisotropy criterion.
\end{remark}
